\section{Discussion and Conclusion}
\label{sec:discussion-conclusion}
\label{sec:lessons}
\label{sec:conclusion}

The experiments point to three findings that a single-axis
leaderboard would miss. First, the FM's value is not a sharper
point forecast but trend, calibration, and robustness upstream and
coverage downstream. Persistence is a strong point-error baseline
and matched-budget precision makes FM-fed and raw-snapshot LLM
variants indistinguishable; yet where persistence is structurally
pinned the FM supplies trend consistency ($0.525$--$0.632$ vs
$0$), calibrated intervals (PI coverage $0.913$, ECE $0.013$), and
$24\%$ lower degradation under stress, and downstream its evidence
supports more targets at undiminished precision and lifts F1.
Second, the LLM failure mode is actionability calibration, not
grounding. Opus~4.7 cites valid fields (grounding $1.000$) but
has $\mathrm{FIR}\!=\!0.437$--$0.448$; Sonnet~4.6 on the same FM
prompts raises F1 to $0.0288$ and reduces FIR to $0.000$. FIR is
therefore a first-class metric alongside grounding and judge
scores. Third, gates are reserve capacity: clean-data effects are
within noise, but under severe masking the data-health gate
suppresses up to $60\%$ of FIR-eligible tickets, so gates need
stress tests and audit logs that support threshold retuning.

Overall, \sysname{} is best read as an auditable
recall--explanation Pareto option for human-in-the-loop DeFi
supervision: structured forecast evidence improves coverage and
rationales, but high-severity recommendations still require
decision-model calibration and intervention-realism measurement.
